\subsection{Analyse de l'énoncé}
Voici ce que je comprends:

\begin{quote}
    Soit un homme qui saute à l'élastique du haut d'un pont de 100 mètres de haut. Il se laisse tomber pour descendre en chute libre pendant 30 mètres, distance à laquelle l'élastique commence à freiner sa descente. C'est donc au bout de 30 mètres qu'il a atteint la vitesse maximale.
    Je dois vérifier si je pourrais moi-même atteindre la vitesse maximale indiquée sur le panneau, qui est de 90 km/h, en sautant depuis ce pont.

    Mais évidemment, il y a de nombreux problèmes dans cet énoncé, comme le fait que le gars sur l'image a la tête en bas quand il arrive 30 mètres plus bas alors qu'il est debout sur le pont, donc la tête en haut. Comment prendre en compte cet élément? Quelle est la taille du gars? S'il mesure 1,80m par exemple, cela veut-il dire que l'élastique ne se tend qu'après 31,80 mètres de chute? À noter que si c'est le cas, et qu'un gars de 2 mètres saute, il atteint les 90 km/h. S'il ne mesure que 1m80, il ne les atteint pas à un pouillème près.
\end{quote}